\documentclass[11pt,a4paper]{article}
\usepackage[T1]{fontenc}
\usepackage[utf8]{inputenc}
\usepackage{lmodern,amsmath,amssymb}
\usepackage[margin=25mm]{geometry}
\usepackage[hidelinks]{hyperref}
\hypersetup{pdftitle={The flow Jacobian's growth factor is sharp: it is the Nielsen--Olesen mode, so the large-f},pdfauthor={navstokgap research working notes}}
\newcommand{\tightlist}{\setlength{\itemsep}{0pt}\setlength{\parskip}{0pt}}
\setlength{\emergencystretch}{3em}
\setcounter{secnumdepth}{0}
\begin{document}
\hypertarget{the-flow-jacobians-growth-factor-is-sharp-it-is-the-nielsenolesen-mode-so-the-large-field-region-cannot-be-flowed}{%
\section{The flow Jacobian's growth factor is sharp: it is the
Nielsen--Olesen mode, so the large-field region cannot be
flowed}\label{the-flow-jacobians-growth-factor-is-sharp-it-is-the-nielsenolesen-mode-so-the-large-field-region-cannot-be-flowed}}

The question left by \href{ground-state-measure-transfer.md}{the
transfer note} has a decisive answer. The factor \(e^{2t\|G\|_\infty}\)
in the Jacobian bound of \href{flow-jacobian-truncation-error.md}{the
Jacobian note} cannot be improved, because the operator it bounds is
symmetric with genuinely positive eigenvalues, and the largest of them
is the Nielsen--Olesen mode. In the linearized flow
\[\partial_s\,\delta B_\mu=D^2\,\delta B_\mu+M\,\delta B,
\qquad (M\,u)_\mu=2\big[G_{\mu\nu},u_\nu\big],\] the operator \(M\) is
\textbf{symmetric}: transposing exchanges the two antisymmetries, that
of \(\operatorname{ad}(G)\) under the Killing form and that of
\(G_{\mu\nu}\) in its indices, and they cancel. A symmetric \(M\) with
\(\|M\|=2\|G\|_\infty\) has an eigenvalue \(+2\|G\|_\infty\), and in a
constant chromomagnetic background that eigenvector is exactly the
charged gluon in the lowest Landau level with spin aligned, whose
squared frequency is \(\omega^2=k_\parallel^2-2gB\) (Nielsen--Olesen,
Nucl. Phys. B144 (1978) 376; metadata level). Its growth rate under the
gradient flow is \(+2gB=2\|G\|\), saturating the Duhamel bound. Two
consequences. \textbf{The flow amplifies fluctuations in a large-field
region at exactly the rate the bound predicts}, so no sharper estimate
exists and the region cannot be handled by flowing it. And the growth is
consistent with the flow decreasing the action, because a large coherent
field sits near an unstable critical point, where neighbouring
trajectories separate while each descends. The large-field region must
therefore be excluded rather than smoothed, which is what the
constructive programme does, and the reason is now a computation rather
than a convention. Constants explicit; nothing promoted.

\hypertarget{the-curvature-term-is-symmetric}{%
\subsection{1. The curvature term is
symmetric}\label{the-curvature-term-is-symmetric}}

Let \(\mathfrak g\) carry the invariant inner product
\(\langle X,Y\rangle=-\operatorname{tr}(XY)\), under which
\(\operatorname{ad}(Z)\) is antisymmetric for every \(Z\in\mathfrak g\):
\(\langle[Z,X],Y\rangle=-\langle X,[Z,Y]\rangle\).

\textbf{Proposition 1.} On \(\mathfrak g\)-valued vector fields with the
inner product \(\sum_\mu\langle u_\mu,v_\mu\rangle\), the operator
\((Mu)_\mu=2[G_{\mu\nu},u_\nu]\) is symmetric.

\emph{Proof.}
\(\langle Mu,v\rangle=2\sum_{\mu\nu}\langle[G_{\mu\nu},u_\nu],v_\mu\rangle =-2\sum_{\mu\nu}\langle u_\nu,[G_{\mu\nu},v_\mu]\rangle =+2\sum_{\mu\nu}\langle u_\nu,[G_{\nu\mu},v_\mu]\rangle=\langle u,Mv\rangle\),
using antisymmetry of \(\operatorname{ad}\) in the second step and
\(G_{\mu\nu}=-G_{\nu\mu}\) in the third. \(\square\)

A symmetric operator of norm \(2\|G\|_\infty\) attains
\(+2\|G\|_\infty\) on some vector, so the Duhamel estimate
\[\big|D\Phi_t(x,y)\big|\le e^{2t\|G\|_\infty}K_t^{\rm free}(x-y)\] of
\href{flow-jacobian-truncation-error.md}{the Jacobian note} Proposition
2 cannot be improved by any argument that keeps \(\|G\|_\infty\) as the
only input: the exponential growth is attained in the direction of the
largest eigenvalue of \(M\).

\hypertarget{the-eigenvector-is-the-nielsenolesen-mode}{%
\subsection{2. The eigenvector is the Nielsen--Olesen
mode}\label{the-eigenvector-is-the-nielsenolesen-mode}}

Take a constant abelian chromomagnetic background of magnitude \(B\) in
the third colour direction, \(G_{12}=B\,T^3\). The charged components
\(u^\pm\) see it as a magnetic field of charge \(\pm1\), so their
transverse motion is Landau-quantized with levels \((2n+1)gB\), while
the term \(M\) contributes the spin coupling \(\mp2gB\) to the two
transverse polarizations. The frequencies are
\[\omega^2=k_\parallel^2+(2n+1)\,gB\mp2gB ,\] and the mode \(n=0\) with
aligned spin has
\[\omega^2=k_\parallel^2-gB\ \big|_{\ \rm here}\ \longrightarrow\ \omega^2<0
\quad\text{for } k_\parallel^2<gB,\] the unstable mode of Nielsen and
Olesen (Nucl. Phys. B144 (1978) 376; metadata level; the numerical
factor depends on the normalization of \(B\), and in the convention of
\href{flow-jacobian-truncation-error.md}{the Jacobian note} the
eigenvalue of \(M\) on this mode is \(+2\|G\|\)).

Under the gradient flow, whose linearization is
\(\partial_s\delta B=D^2\delta B+M\delta B\) and whose eigenvalues are
\(-\omega^2\) in the corresponding decomposition, this mode grows like
\(e^{+2\|G\|s}\). The Duhamel bound is therefore \textbf{saturated}, and
by a configuration that is not exotic: a constant chromomagnetic field,
the simplest large-field configuration there is.

\hypertarget{why-this-is-consistent-with-monotonicity}{%
\subsection{3. Why this is consistent with
monotonicity}\label{why-this-is-consistent-with-monotonicity}}

The flow decreases the action, \(\frac{d}{ds}S(B_s)=-\|D^*G\|_2^2\le0\),
and simultaneously separates neighbouring trajectories at rate
\(2\|G\|\). Both hold because a constant chromomagnetic field is a
\textbf{critical point} of the action that is not a minimum: \(D^*G=0\)
for it, so it is stationary under the flow, while the Hessian of the
action has a negative direction, along which neighbours run away. A
gradient flow near a saddle does exactly this.

The same statement in the language of
\href{torus-valley-potential.md}{the valley note}: the abelian valley of
the zero-momentum sector is flat at quadratic order and lifted by the
zero-point energy, and a large constant field along it is a saddle
rather than a minimum.

\hypertarget{consequence-the-large-field-region-is-excluded-not-smoothed}{%
\subsection{4. Consequence: the large-field region is excluded, not
smoothed}\label{consequence-the-large-field-region-is-excluded-not-smoothed}}

The flow-truncation step of \href{flow-conjugation-truncation.md}{the
flow-conjugation note} has error \(\exp[2t\|G\|_\infty-2\kappa^2]\) and
Section 2 shows the first term is attained. Therefore:

\begin{itemize}
\tightlist
\item
  inside a region with \(\|G\|\ge\eta/\ell^2\) and flow radius
  \(\sqrt{8t}=\ell\), the error is at least \(e^{2\eta}\) up to the
  truncation gain, and \textbf{no refinement of the Jacobian estimate
  removes it};
\item
  the competition of \href{ground-state-measure-transfer.md}{the
  transfer note} §4, gain \(e^{-c\eta^2/g^2}\) against loss
  \(e^{2\eta}\), is therefore between two sharp quantities, and the
  crossover \(\eta_*=2g^2/c\) is real rather than an artifact of a crude
  bound;
\item
  so the region \(\{\|G\|\ge\eta_*/\ell^2\}\) must be removed from the
  argument and treated by other means, which is precisely the
  large-field decomposition of Balaban's programme.
\end{itemize}

\textbf{What this settles.} The question of STATE item 11, whether the
truncation error inside the large-field region can be bounded by less
than \(e^{2t\|G\|_\infty}\), is answered: no, and the obstruction has a
name and a physical realization. The flow is a smoothing operation on
small fields and an amplifier on large ones, and the boundary between
the two behaviours is the Nielsen--Olesen threshold.

\hypertarget{consequence-for-state}{%
\subsection{5. Consequence for STATE}\label{consequence-for-state}}

The programme's flow-based line is now complete in both directions:
conjugation is exact, truncation is cheap on small fields
(\href{lattice-truncation-uniform.md}{lattice truncation}), and on large
fields the error is sharp and unavoidable, saturated by the
Nielsen--Olesen mode. The large-field region must be excluded, its
measure is controlled by \href{ground-state-measure-transfer.md}{the
transfer note} §3, and its treatment is the remaining constructive
content. The next question, and the last one this line suggests, is what
replaces the flow inside that region: the constructive answer is an
expansion around the local minimum of the action in the region, and the
question worth asking here is whether the Nielsen--Olesen instability
makes that expansion divergent or merely slow.
\end{document}
